\def\title{Determining k Clusters in k-Means Clustering with the CRAB Algorithm}
\def\author{Jasmine Kristine Soriano Cabrera}
\def\degreemonth{June}
\def\degreeyear{2026}
\def\degree{Master of Science in Statistics}
\def\numberofmembers{3}
   \def\chair{Kelly Bodwin, Ph.D. \linebreak Professor of Statistics}
   \def\othermemberA{Allison Theobold, Ph.D. \linebreak Professor of Statistics}
   \def\othermemberB{Julia Schedler, Ph.D. \linebreak Professor of Statistics}
\def\keywords{Clustering, Data Science, k-Means.} % Make sure to end this in a . (i.e. Some, Thing, Interesting.)

\def\abstract{%
Unsupervised clustering often faces the challenge of determining the correct number 
of clusters in the absence of a true target variable. Traditional methods such
as the Elbow Method and the Silhouette Score can produce ambiguous results and rely on 
assumptions about cluster shape or separation. To address this, we created the 
Clustering Rivals and Buddies (CRAB) algorithm which evaluates clusters based on 
stability across multiple subsamples. CRAB uses pairwise classifications to identify 
points that consistently group together called "Buddies" and points that remain 
separated caled "Rivals". Applied with k-Means, CRAB accurately recovers underlying 
cluster structures in both spherical and non-spherical datasets, even under varying 
sample sizes and variances. An accompanying R package implements CRAB with a tidy 
workflow and easy to follow syntax, making it accessible for research and educational 
use. By emphasizing stability, CRAB offers a robust alternative for selecting cluster 
numbers in datasets where traditional methods may fall short.

}

\def\tocpagelegnth{1} % Number of pages your Table of Contents is (used to set correct page numbers for List of Tables and List of Figures)
\def\totflag{1} % 1 if you have at least 1 table, otherwise 0 (adds Table of Tables page)
\def\tofflag{1} % 1 if you have at least 1 figure, otherwise 0 (adds Table of Figures page)
